% ------------------------------------------------------------------------------
% Conclusão
% ------------------------------------------------------------------------------


\chapter{Conclusão}
\label{chap_conclusao}


O estudo da queda livre, apesar de sua simplicidade aparente, mostrou-se um excelente laboratório para a aplicação e análise de métodos numéricos, modelagem matemática e validação de soluções analíticas. As simulações computacionais permitiram comparar regimes físicos distintos — sem resistência do ar, com resistência linear e quadrática — e avaliar o impacto dos parâmetros físicos e numéricos na precisão dos resultados.

Os resultados evidenciaram que a escolha do passo de tempo ($\Delta t$) é determinante para a acurácia dos métodos numéricos, especialmente em modelos com resistência do ar, que exigem discretizações mais refinadas para garantir resultados confiáveis. A comparação entre soluções numéricas e analíticas demonstrou a convergência dos métodos implementados, além de revelar limitações e potencialidades de cada abordagem.

Além de aprofundar a compreensão dos fenômenos físicos envolvidos, este trabalho reforça a importância da modelagem matemática e da análise computacional na solução de problemas reais, destacando a necessidade de escolhas criteriosas de parâmetros e validação rigorosa dos resultados.

Como perspectivas futuras, recomenda-se a investigação de outros métodos numéricos, a análise de cenários com diferentes condições iniciais e a aplicação dos conceitos desenvolvidos em contextos experimentais ou interdisciplinares. Espera-se que este estudo contribua para a formação científica e tecnológica, servindo de base para novas pesquisas e aplicações didáticas.
 

\section*{Comparação dos Perfis de Movimento}

Inicialmente, foi realizado o comparativo entre as três simulações principais: queda livre sem resistência do ar, com resistência linear e com resistência quadrática. Os gráficos a seguir apresentam a evolução de posição, velocidade e aceleração ao longo do tempo para cada modelo físico:

\begin{figure}[H]
	\centering
	\begin{subfigure}[b]{0.32\textwidth}
		\centering
		\includegraphics[width=\textwidth]{./figuras/graficos/comparativo/comparativo_posicao.png}
		\caption{(a) Posição}
		\label{fig:comparativo_posicao}
	\end{subfigure}
	\hfill
	\begin{subfigure}[b]{0.32\textwidth}
		\centering
		\includegraphics[width=\textwidth]{./figuras/graficos/comparativo/comparativo_velocidade.png}
		\caption{(b) Velocidade}
		\label{fig:comparativo_velocidade}
	\end{subfigure}
	\hfill
	\begin{subfigure}[b]{0.32\textwidth}
		\centering
		\includegraphics[width=\textwidth]{./figuras/graficos/comparativo/comparativo_aceleracao.png}
		\caption{(c) Aceleração}
		\label{fig:comparativo_aceleracao}
	\end{subfigure}
	\caption{Comparação dos perfis de (a) posição, (b) velocidade e (c) aceleração ao longo do tempo para os diferentes modelos.}
	\label{fig:comparativo_abc}
\end{figure}

Esses resultados gráficos permitem visualizar de forma clara as diferenças dinâmicas entre os modelos, evidenciando o impacto da resistência do ar sobre o movimento e a importância da modelagem adequada para cada regime físico.

\section*{Síntese dos Resultados de Acurácia}

Na sequência, foi realizada a análise de acurácia dos métodos numéricos empregados. O gráfico a seguir ilustra o erro relativo médio para cada modelo físico e método numérico, em função do passo de tempo ($\Delta t$):

\begin{figure}[H]
	\centering
	\includegraphics[width=0.8\textwidth]{./figuras/graficos/acuracia/acuracia_comparativa.png}
	\caption{Comparação do erro relativo médio entre os métodos numéricos e modelos físicos, em função de $\Delta t$.}
	\label{fig:acuracia_comparativa}
\end{figure}

Para consolidar a análise, apresenta-se a seguir a tabela comparativa de acurácia obtida nos experimentos:

\begin{center}
	\input{./tabelas/output/acuracia/tabela_acuracia.tex}
\end{center}

Os resultados evidenciam que a escolha do passo de tempo ($\Delta t$) influencia diretamente a precisão dos métodos numéricos, sendo possível observar a convergência dos resultados para valores menores de $\Delta t$. Modelos com resistência do ar apresentam maior sensibilidade ao passo de tempo, exigindo discretizações mais finas para garantir acurácia adequada.

Essas análises reforçam a importância da escolha criteriosa dos parâmetros numéricos e validam a abordagem computacional adotada neste trabalho.

